%TeX root = ../main.tex

\chapter{Abstract Integration}

\defn{Sigma Algebra}{Let $\bbM \subseteq \mathscr P(X)$ for a given set $X$ that we call the \emph{Sample space}. We say $\bbM$ is a sigma algebra if \begin{enumerate}
    \item (Closed under countable union) If $E_1$,$E_2$, $E_3 \cdots$ are in $\bbM$, then $\cup_i E_i \in \bbM$ 
    \item $\varnothing$ and $X$ are in $\bbM$
    \item (Closed under complementation) If $E\in \bbM$ then $E^c \in \bbM$
\end{enumerate}}
\defn{Sigma Algebra generated by a set}{Let $S$ be a collection of subsets of $X$. Then $\sigma(S)$ is defined to be the smallest sigma algebra that contains $S$ as a subcollection.}
\defn{Borel Sigma Algebra}{Let $X,\mathscr T$ be a topological space. The borel sigma algebra $\mathscr B(X)$ is the sigma algebra generated by all the open sets in $X$. i.e. $\sigma(\mathscr T)=\mathscr B(X) $}
\defn{Measurable Function}{A function $f:X, \bbM_X\to Y,\bbM_Y$ is called \emph{measurable} if for every $S\in \bbM_Y$, $f^{-1}(S) \in \bbM_X$
\\
In context of real or complex valued functions, we take our $\bbM_Y$ to be the corresponding borel sigma algebra. }
\defn{Measure}{Let ($X,\bbM_X$) be a measurable space. A set function $\mu:\bbM_X\to [0,\infty]$ is called a \emph{positive measure} if \begin{enumerate}
    \item $\exists S_0 \in \bbM_X$ such that $\mu(S_0)<\infty$
    \item If $\{E_i\}_{i=1}^{\infty}$ is a collection of pairwise disjoint sets in $\bbM_X$, then $\mu(\cup_i E_i)=\sum_{i=1}^{\infty} \mu(E_i)$ (\emph{\textbf{Countable Additivity}})
\end{enumerate}}

\thmp{}{Let $f:(X,\bbM_X)\to (Y,\bbM_Y=\sigma(\mathcal F))$. $f$ is measurable if and only if for every $F\in \mathcal F$, $f^{-1}(F)\in \bbM_X$}{$\implies$) If $f:X\to Y$ is measurable, then for every $F\in \bbM_Y$, we have $f^{-1}(F)\in \bbM_X$, which means $\forall F \in \cal F$ we have $f^{-1}(F)\in \bbM_X$
\\ $\impliedby$) Suppose for every $F\in \cal F $ we have $f^{-1}(F)\in\bbM_X$. Let $B=\{E\in \bbM_Y: f^{-1}(E)\in \bbM_X\}$. Clearly by hypothesis, $\cal F \subset B$. Clearly $\varnothing$ and $Y$ are in $B$. Let $E\in B$, i.e. $f^{-1}(E)\in \bbM_X$. Since $(f^{-1}(E))^c \in \bbM_X$, we have $f^{-1}(E^c)\in \bbM_X$ or $E^c\in B$. Hence, $B$ is closed under complements. Let $E_1,E_2 \cdots E_n \cdots$ be a collection of sets in $B$. $f^{-1}(\cup_i E_i)=\cup_i f^{-1}(E_i) \in \bbM_X$ which means $\cup_i E_i$ is in $B$. Hence $B$ is a sigma algebra, hence $\sigma(\cal F)=\bbM_Y \subseteq B \subseteq \bbM_Y$ which means $\bbM_Y=B$ (\textbf{This is an example of a technique called \emph{Method of Good Sets}}). }

\thmp{\label{Continuous composition}}{Let $Y$ and $Z$ be topological spaces, Let $g:Y\to Z$ be a continuous mapping. If $(X,\Sigma(X))$ is a measurable space, and $f:X\to Y$ is measurable (in the Borel Sense), then $g\circ f:X\to Z$ is also measurable.}{If $g:Y\to Z$ is continuous, then $g^{-1}(V)$ is open in $Y$, for every $V$ open in $Z$. Consider $(g\circ f)^{-1}(V)=f^{-1}(g^{-1}(V))$. This is clearly in $\Sigma(X)$ since $g^{-1}(V)$ is open and $f$ is measurable. Hence $g\circ f$ is measurable. 

% https://q.uiver.app/#q=WzAsMyxbMiwwLCJZIl0sWzQsMCwiWiJdLFswLDAsIihYLFxcU2lnbWEoWCkpIl0sWzAsMSwiZyAoXFx0ZXh0e2NvbnRpbnVvdXN9KSIsMl0sWzIsMCwiZiAoXFx0ZXh0e21lYXN1cmFibGV9KSIsMl0sWzIsMSwiZ1xcY2lyYyBmIiwxLHsibGFiZWxfcG9zaXRpb24iOjQwLCJvZmZzZXQiOjIsImN1cnZlIjo1fV1d
\[\begin{tikzcd}[ampersand replacement=\&]
	{(X,\Sigma(X))} \&\& Y \&\& Z
	\arrow["{f (\text{measurable})}"', from=1-1, to=1-3]
	\arrow["{g\circ f \text{(measurable)}}"{description, pos=0.4}, shift right=2, curve={height=30pt}, from=1-1, to=1-5]
	\arrow["{g (\text{continuous})}"', from=1-3, to=1-5]
\end{tikzcd}\]

}

\thmp{}{Let $u:X\to \bbR$ and $v:X\to \bbR$ be measurable, real mappings. Let $\varphi: \bbR^2 \to Y$ be a continuous map from the plane to the topological space $Y$. Let $h:X\to Y$ be $h(x)=\varphi(u(x),v(x))$
Then $h$ is a measurable function from $X$ to $Y$. % https://q.uiver.app/#q=WzAsNSxbMCwxLCJYIl0sWzQsMSwiWSJdLFsyLDAsIlxcbWF0aGJiICBSIl0sWzIsMiwiXFxtYXRoYmIgUiJdLFsyLDEsIlxcbWF0aGJiIFJeMiJdLFswLDIsInUiXSxbMCwzLCJ2IiwyXSxbNCwxLCJcXHZhcnBoaSIsMl0sWzAsNCwiKHUoeCksdih4KSkiLDJdXQ==
\[\begin{tikzcd}[ampersand replacement=\&]
	\&\& {\mathbb  R} \\
	X \&\& {\mathbb R^2} \&\& Y \\
	\&\& {\mathbb R}
	\arrow["u", from=2-1, to=1-3]
	\arrow["{(u(x),v(x))}"', from=2-1, to=2-3]
	\arrow["v"', from=2-1, to=3-3]
	\arrow["\varphi"', from=2-3, to=2-5]
\end{tikzcd}\]

}{$h(x)=\varphi(u(x),v(x))$. Let $V$ be open in $Y$. $\varphi^{-1}(V) $ is open in $\bbR^2$ so it is the countable union of open sets of the form $A_i \times B_i$ where $A_i$ and $B_i$ are each open in $\bbR$. $\varphi^{-1}(V)=\cup_i (A_i \times B_i)$. $(u,v)^{-1}(\varphi^{-1}(V))=(u,v)^{-1}(\cup_i A_i \times B_i)=\cup_i (u,v)^{-1}(A_i \times B_i)$. $(u,v)^{-1}(A_i \times B_i):=\{x\in X: u(x)\in A_i, v(x)\in B_i \}=u^{-1}(A_i)\cap v^{-1}(B_i) $. Therefore, $h^{-1}(V)=\cup_i (u^{-1}(A_i)\cap v^{-1}(B_i))$ which is a countable union of measurable sets, which is measurable. Hence, $h$ is a measurable mapping.
\\ \emph{Note:} It was enough to show the map $X\xrightarrow[]{(u(x),v(x))} \bbR^2$ is measurable, from theorem \ref{Continuous composition}
}
\thmp{}{If $u$ and $v$ are real valued, measurable functions from a measurable space $X$, then $f(x)=u(x)+iv(x)$ is a complex valued measurable function.}{$X\xrightarrow[\text{measurable}]{(u(x),v(x))} \bbR^2 \xrightarrow[]{\text{identity}}\bbC $ composition gives us $f(x)=u(x)+iv(x)$. Since $(u(x),v(x))$ is measurable and identity is continuous, we get that $f$ is a measurable map. Therefore, it is enough to check that the real and imaginary parts of a complex function are measurable as real valued functions.  }
\corp{If $f:X\to \bbC$ is a complex valued measurable function, then $Re(f)$ and $im(f)$ are real valued measurable functions. Moreover, $|f|$, the complex modulus, is also measurable. }{$X\xrightarrow[]{f}\bbC \xrightarrow[]{Re(x)/Im(x)/|x|} \bbR$. Since $f$ is measurable and $Re(x)/Im(x)/ |x|$ are continuous we have that the compositions $Re(f)$, $|f|$ or $Im(f)$ are themselves measurable. }
\rmk{If $f:X\to \bbR\cup\{\pm \infty\}$ is given, then we call $f$ measurable if for every $x\in \bbR$ we have $f^{-1}(a,\infty] \in \bbM_X$ (Think of the generating sets of the order topology of $\bbR \cup \{\pm \infty\}$) }
\lemp{(Positive and Negative Parts of a Real Valued Function)}{Every function $f:X\to \bbR\cup\{\pm \infty\}$ can be split into a sum of two measurable, positive valued functions.}{$f:X\to \bbR\cup\{\pm \infty\}$. Let $f^+(x):=\max\{f(x),0\}$ and $f^{-}:=\max\{-f(x),0\}$. Then it is clear that $f=f^+-f^-$ and $|f|=f^++f^-$.}
\thmp{}{If $f$ and $g$ are measurable, extended real valued mappings, then $f+g$ and $f\cdot g$ are also measurable, extended real valued mappings.}{$(f+g)^{-1}(a,\infty]:=\{x \in X: a-g(x)<f(x)\}=\{x\in X:\exists r\in \bbQ \text{ such that } a-g(x)<r<f(x) \}=$ $\{x\in X: \exists \text{ such that }a-g(x)<r \text{ and } r<f(x) \}=\cup_{r \in \bbQ}\{x \in X:a-g(x)<r \textbf{ and }r<f(x)\}=$ $\cup_{r\in \bbQ} (\{x\in X: a-r<g(x)\}\cap \{x\in X: r<f(x)\})$ which is eventually a measurable set. Hence $f+g$ is measurable.
\\\\ Let $k \in 2\bbN$. Let $g=f^k$. $g^{-1}(a,\infty]=\{x\in X: a<f^k\}=\{x\in X: a^{1/k}<f \textbf{ and }f(x)<-a^{1/k} \}$ which is a measurable set (Regardless of whether $a$ is positive or negative). If $k \in 2\bbN +1$ then $g^{-1}(a,\infty]=\{x\in X: a<f(x)^k\}=\{x:a^{1/k}<f(x) \}$. If $a<0$, $a=-t$ and we have $\{x\in X: -t<f(x)^{k}\}=\{x\in X: -t^{1/k}<f(x) \}$. All these sets are eventually measurable. Hence, $f^k$ is measurable.
\\\\ $1/4((f+g)^2-(f-g)^2)=f\cdot g$. Since addition of functions and raising functions to powers is measurable, we have $f\cdot g$ is also measurable.   }
\lemp{}{$E\subset X$ is measurable as a set if and only if $\chi_E$ is measurable as a function from $X \xrightarrow[]{\chi_E} \bbR$}{$\implies$) Let $E$ be a measurable set. $\chi_E(x):=\begin{cases}
    1 \ \text{ if }x \in E \\
    0 \ \text{ otherwise}
\end{cases}$ Look at $\chi_E^{-1}(a,\infty]$. If $1\leq  a$, then $\chi_E^{-1}(a,\infty]$ is empty, which is measurable. If $a<1$, then $\chi_E^{-1}(a,\infty]=\{x\in X: \chi_E(x)\in (a,\infty]\}=E$ which is measurable. If $a<0$, then the inverse image is $X$ which is measurable. Hence, $\chi_E$ is a measurable function.
\\\\ $\impliedby$) Let $\chi_E:X\to \bbR$ be measurable. Note that $\{1\}$ singleton is in the borel sigma algebra of $\bbR$. Hence, we see $\chi_E^{-1}(\{1\})=E$ is necessarily measurable. Therefore $E$ is measurable.  }

\lemp{}{If $f:X\to \bbC$ is measurable, then there exists $\alpha:X\to \bbC$ that is measurable such that $|\alpha|=1$ and $f=|f|\alpha$.}{Let $\alpha(x):=\begin{cases}
    1 \ \text{ if } f(x)=0 \\
    f/|f| \text{ if } f(x)\neq 0
\end{cases}$ 
\\ Look at $\alpha^{-1}(t,\infty]=\{x\in X: \alpha(x)>t\}$ We can disjointify $X$ into $E:=\{x \in X: f(x)=0\}$ and $E^c:=\{x\in X: f(x)\neq 0\}$. We can then write $\alpha^{-1}(t,\infty]=\alpha^{-1}(t,\infty]\cap(X)=$ $(\alpha^{-1}(t,\infty]\cap E)\cup(\alpha^{-1}(t,\infty]\cap E^c)$ }
